\documentclass[../../../include/open-logic-section]{subfiles}

\begin{document}

\olfileid{sth}{replacement}{refproofs}
\olsection{Lampiran: Asil-asil babagan Panggantian}

Ing bagean iki, kita bakal mbuktekake Refleksi
ing $\ZF$. Kita uga bakal mbuktekake salah siji
pangerten yen Refleksi ekivalen karo Panggantian.
Sabanjure kita mbuktekake akibat kang narik
kawigaten babagan kakuwatan Refleksi/Panggantian.
\emph{Pènget: bagean iki kalebu matematika
kang paling maju ing buku iki.}

Kita miwiti saka lema kang, supaya ringkes,
nganggo notasi \emph{garis ndhuwur} kanggo
urutan variabel utawa objek. Dadi,
``$\overline{a}_k$'' nyingkat
``$a_{k_1}$, \dots, $a_{k_n}$'',
lan $n$ ditemtokake dening konteks.

\begin{lem}\ollabel{lemreflection}
Kanggo saben $1 \leq i \leq k$, tetepna
$\phi_i(\overline{v}_i, x)$ minangka formula.
Mula kanggo saben $\alpha$ ana sawijining
$\beta > \alpha$ kang nyukupi pratelan iki
kanggo saben $\overline{a}_1, \ldots,
\overline{a}_k \in V_\beta$ lan
$1 \leq i \leq k$:
\[
	\exists x\phi_i(\overline{a}_i, x) \rightarrow (\exists x \in V_\beta) \phi_i(\overline{a}_i, x)
\]
\end{lem}

\begin{proof}
Kita netepake term $\mu$ mangkene:
$\mu(\overline{a}_1, \ldots,
\overline{a}_k)$ iku tataran paling awal, $V$,
kang nyukupi kabeh implikasi iki kanggo
$1 \leq i \leq k$:
\[
\exists x\phi_i(\overline{a}_i, x) \rightarrow (\exists x \in V) \phi_i(\overline{a}_i, x)
\]
Gampang dipriksa yen
$\mu(\overline{a}_1, \ldots, \overline{a}_k)$
ana kanggo saben $\overline{a}_1, \ldots,
\overline{a}_k$. Saiki, kanthi Panggantian lan
teorema rekursi kita, tetepna:
\begin{align*}
	S_0 & = V_{\alpha+1}\\
	S_{n+1} & = S_n \cup \bigcup
	\Setabs{\mu(\overline{a}_1, \ldots, \overline{a}_k)}
	{\overline{a}_1, \ldots, \overline{a}_k \in S_n} \\
	S &= \bigcup_{m < \omega} S_m.
\end{align*}
Saben $S_n$, lan mula $S$ dhewe, iku tataran
sawise $V_\alpha$. Saiki tetepna
$\overline{a}_1$, \dots,~$\overline{a}_k \in S$;
mula ana sawijining $n < \omega$ kang nyukupi
$\overline{a}_1$, \dots,
$\overline{a}_k \in S_n$. Tetepna sawijining
$1 \leq i \leq k$, lan upamane
$\exists x \phi_i(\overline{a}_i,x)$.
Mula $(\exists x \in \mu(\overline{a}_1,
\ldots, \overline{a}_k))\phi_i(\overline{a}_i, x)$
miturut konstruksi, saengga
$(\exists x \in S_{n+1})\phi_i(\overline{a}_i, x)$,
lan banjur
$(\exists x \in S)\phi_i(\overline{a}_i, x)$.
Dadi $S$ iku $V_\beta$ kang digoleki.
\end{proof}
\noindent 
Saiki kita bisa mbuktekake
\olref[sth][replacement][ref]{reflectionschema}
kanthi cukup langsung:

\begin{proof}[Bukti]
Tetepna $\alpha$. Tanpa ngurangi keumuman,
kita bisa nganggep yen panyambung ing $\phi$
mung $\exists$, $\lnot$, lan $\land$ (amarga
kabeh iki cukup kanggo nyatakake formula).
Tetepna $\psi_1, \ldots, \psi_k$ minangka
dhaptar subformula $\phi$ miturut
kerumitane, saengga $\psi_k = \phi$.
Miturut \olref{lemreflection}, ana
$\beta > \alpha$ kang nyukupi pratelan iki
kanggo saben $\overline{a}_i \in V_\beta$
lan $1 \leq i \leq k$:
\begin{align}\label{reflectionnicelybehaved}
	\exists x\psi_i(\overline{a}_i, x) \rightarrow 
	(\exists x \in V_\beta) \psi_i(\overline{a}_i, x)\tag{*}
\end{align}
Kanthi induksi ing kerumitan $\psi_i$,
kita bakal nuduhake yen
$\psi_i(\overline{a}_i) \leftrightarrow
\psi_i^{V_\beta}(\overline{a}_i)$ kanggo saben
$\overline{a}_i \in V_\beta$. Yen $\psi_i$
atomik, iki cetha. Bikondisional iki uga
njamin yen $\psi_i$ arupa negasi utawa
konjungsi saka subformula kang wis nyukupi
sipat iki, mula $\psi_i$ uga nyukupine.
Dadi mung kasus kuantifikasi kang perlu
dipriksa. Tetepna $\overline{a}_i \in V_\beta$;
mula:
\begin{align*}
	(\exists x \psi_i(\overline{a}_i, x))^{V_\beta}
	&\text{ yen lan mung yen }
	(\exists x \in V_\beta)\psi_i^{V_\beta}(\overline{a}_i, x)
	&&\text{miturut definisi}\\
	&\text{ yen lan mung yen }
	(\exists x \in V_\beta)\psi_i(\overline{a}_i,  x)
	&&\text{miturut pangira-ira}\\
	&\text{ yen lan mung yen }
	\exists x \psi_i(\overline{a}_i, x)
	&&\text{miturut \eqref{reflectionnicelybehaved}}
\end{align*}
Iki ngrampungake induksi; asil kasebut
ndherek amarga $\psi_k = \phi$.
\end{proof}

Kita wis mbuktekake Refleksi ing $\ZF$.
Bukti iki umume manut \citet{Montague1961}.
Saiki kita arep mbuktekake ing $\Z$ yen
Refleksi ngakibatake Panggantian. Bukti
iki manut \citet{Levy1960}, nanging luwih
disederhanakake.

Amarga kita makarya ing $\Z$, Refleksi
ora bisa ditulis persis kaya ing ndhuwur.
Rumusan mau nganggo notasi ``$V_\alpha$'',
kang ora bisa ditetepake ing $\Z$ (delengen
\olref[sth][spine][zf]{sec}). Mula kita
ngusulake rumusan Panggantian kang katon
luwih ringkih mangkene:

\begin{defish}
\emph{Refleksi Lemah.} Kanggo saben formula
$\phi$, ana himpunan transitif $S$ kang ngemot
$0$, $1$, lan saben parameter ing $\phi$
minangka !!{element}s saka $S$, sarta
$(\forall \overline{x} \in S)(\phi \liff
\phi^S)$.
\end{defish}

Kanggo nggunakake iki kanggo mbuktekake
Panggantian, kita bakal manut
\citet[bagéan kapisan Teorema 2]{Levy1960}
lan nuduhake yen rong formula bisa
``direfleksikake'' bebarengan:

\begin{lem}[ing $\Z + \text{Refleksi Lemah}$.]\ollabel{lem:reflect}
Kanggo saben formula $\psi, \chi$, ana
himpunan transitif $S$ kang ngemot $0$ lan
$1$ (lan saben parameter ing formula-formula
mau) minangka !!{element}s saka $S$, sarta
$(\forall \overline{x} \in S)((\psi \liff
\psi^S) \land (\chi \liff \chi^S))$.
\end{lem}

\begin{proof}
Tetepna $\phi$ minangka formula
$(z = 0 \land \psi) \lor (z = 1 \land \chi)$,
kanthi variabel pananda anyar.

Ing kene kita nganggo cekakan: ``$z = 0$''
kudu ditulis kanthi lengkap minangka
``$\forall t\, t \notin z$'', lan ``$z =1$''
minangka ``$\forall s(s \in z \liff
\forall t\, t \notin s)$''. Nanging amarga
$0, 1 \in S$ lan $S$ transitif, formula-formula
iki \emph{absolut} tumrap $S$; tegese,
formula mau lumaku kanggo objek kang padha
sanadyan kuantore diwatesi marang
$S$.\footnote{Luwih formal, yen $\xi$ salah
siji saka formula mau, mula
$\xi(z) \liff \xi^S(z)$.}

Miturut Refleksi Lemah, ana $S$ kang trep
lan nyukupi:
\begin{align*}
	(\forall z, \overline{x} \in S)(&\phi \liff \phi^S)\\
	\text{yaiku }(\forall z, \overline{x} \in S)(&((z = 0 \land \psi) \lor (z = 1 \land \chi)) \liff {}\\
	&\phantom{(}((z = 0 \land \psi) \lor (z = 1 \land \chi))^S)\\
	\text{yaiku }(\forall z, \overline{x} \in S)(&((z = 0 \land \psi) \lor (z = 1 \land \chi))\liff {}\\
	&\phantom{(}((z = 0 \land \psi^S) \lor (z = 1 \land \chi^S)))\\
	\text{yaiku }(\forall \overline{x} \in S)(&(\psi \liff \psi^S) \land (\chi \liff \chi^S))
\end{align*}
Pratelan kapindho ngakibatake kang katelu
amarga ``$z = 0$'' lan ``$z=1$'' absolut
tumrap $S$; pratelan kapapat ndherek amarga
$0 \neq 1$.
\end{proof}\noindent Saiki kita bisa
mbuktekake Panggantian kanthi manut lan
nyederhanakake \citet[Teorema 6]{Levy1960}:

\begin{thm}[ing $\Z$ + Refleksi Lemah]\label{thm:replacement}
Kanggo saben formula $\phi(v,w)$ lan saben
$A$, yen $(\forall x \in A)\lexists![y][\phi(x,y)]$,
mula $\Setabs{y}{(\exists x \in A)\phi(x,y)}$ ana.
\end{thm}

\begin{proof}
Tetepna $A$ kang nyukupi
$(\forall x \in A)\lexists![y][\phi(x,y)]$,
lan tetepna formula-formula iki:
\begin{align*}
	\psi &\text{ iku } (\phi(x, z) \land A = A)\\
	\chi &\text{ iku } \lexists[y][\phi(x, y)]
\end{align*}
Kanthi \olref{lem:reflect}, amarga $A$
parameter ing $\psi$, ana himpunan
transitif~$S$ kang nyukupi $0, 1, A \in S$
(bebarengan karo saben parameter liyane),
lan uga:
\[
	(\forall x,z \in S)((\psi \liff \psi^S) \land (\chi \liff \chi^S))
\]
Dadi, mligine:
\begin{align*}
	(\forall  x, z \in S)(&\phi(x, z) \liff \phi^S(x, z))\\
	(\forall x \in S)(&\exists y\phi(x, y) \liff (\exists y \in S)\phi^S(x, y)) 
\end{align*}
Kanthi nggabungake iki lan ngelingi yen
$A \subseteq S$ amarga $A \in S$ lan $S$
transitif:
\begin{align*}
	(\forall x \in A)(&\exists y\phi(x, y) \liff (\exists y \in S)\phi(x, y))
\end{align*}
Saiki $(\forall x \in A)(\lexists![y \in S])\phi(x, y)$,
amarga $(\forall x \in A)\lexists![y][\phi(x, y)]$.
Miturut Pamisahan, kita oleh
$\Setabs{y \in S}{(\exists x \in A) \phi(x, y)}
= \Setabs{y}{(\exists x \in A) \phi(x, y)}$.
\end{proof}

\end{document}
